\documentclass[twocolumn]{article}
\usepackage[utf8]{inputenc}
\usepackage{amsmath}
\usepackage{amssymb}
\usepackage{graphicx}
\usepackage{booktabs}
\usepackage{hyperref}
\usepackage{listings}
\usepackage{color}

\title{Optimizing Large Generative AI Models for Local CPU Execution via Virtual Address Space Mapping and BFloat16 Numerical Compression:\\\large Empirical Analysis of Zero-RAM Weight Streaming and Multithreaded Throttling on 1.7B Parameter Qwen3-TTS}
\author{andrea lew}
\date{May 2026}

\definecolor{codegray}{rgb}{0.95,0.95,0.95}
\definecolor{codeblue}{rgb}{0.1,0.2,0.4}

\lstset{
    backgroundcolor=\color{codegray},
    basicstyle=\small\ttfamily,
    keywordstyle=\color{codeblue}\bfseries,
    frame=single,
    breaklines=true,
    language=Python
}

\begin{document}

\maketitle

\begin{abstract}
This research presents and validates an ultra-efficient inference framework designed to bypass the high-end GPU and massive system RAM requirements that block the deployment of Large Generative AI models on low-spec client machines and legacy CPU-only devices. Using the 1.7B parameter multilingual Text-to-Speech model Qwen3-TTS as an empirical target, we developed an integrated local execution engine. This engine combines kernel-level virtual address space mapping ($\mathtt{low\_cpu\_mem\_usage=True}$), numerical float compression via 16-bit Brain Float ($\mathtt{bfloat16}$), and thread limit throttling to control context-switching overheads.

Empirical benchmarks reveal that our optimization framework completely resolves the structural bottlenecks of default Eager Loading. Under default Float32 eager loading, initializing the model (~3.8 GB) incurs a 1,727.28 MB heap RAM overhead and a 25.76-second initialization delay. In contrast, under our optimized framework, model loading time is cut to 8.28 seconds (a 3.1x speedup), while active physical RAM allocation spikes are compressed to 0 MB or below (registering an active reduction of $-428.54$ MB due to lazy initialization and garbage collection).

Additionally, we engineered a desktop controller ($\mathtt{gui\_runner.py}$) incorporating async background workers, real-time RSS physical memory tracking, and native winsound audio playback to prevent UI thread locks. Finally, we derive the information-theoretic alignment between virtual neural weight streaming and a 15-bit phonological bit-packing compression algorithm optimized for resource-constrained Korean language processing.
\end{abstract}

\section{Introduction}
Modern generative AI architectures have grown exponentially in parameter size, significantly enhancing representational capacity but at the expense of massive hardware requirements, especially dedicated GPU bandwidth and physical system RAM. Consequently, deploying state-of-the-art neural networks locally on legacy client nodes (e.g., 8GB RAM laptops) has become highly impractical.

This physical resource bottleneck is conceptually equivalent to the severe memory allocation problems solved in the early computer science and COBOL era using virtual memory mapping ($\mathtt{mmap}$) and hardware-level paging. Because generative models evaluate sequential layers during a feed-forward pass, maintaining the weights of all layers in active RAM simultaneously is highly inefficient. At any moment, only the current layer $L_i$ is actively computed; the weights of all other layers $\{L_j \,|\, j \neq i\}$ are static and can remain offline.

This paper proposes and validates a \textbf{Zero-RAM Weight Streaming Inference Framework} by combining virtual address space mapping with 16-bit Brain Float compression. Instead of eagerly parsing the entire network into the heap, parameters are mapped directly into the operating system's page cache and loaded into physical RAM on-demand (lazy page-in) during sequential forward passes.

\section{Theoretical Background and Mathematical Formalization}
\subsection{Mathematical Modeling of Lazy Virtual Address Space Mapping}
Under standard Eager Loading, a model parameter file $F \in \mathbb{R}^{M}$ ($M$ is the total model size in bytes) is read entirely into the heap memory space $\Omega_{\text{heap}}$. The resident set size (RSS) memory consumption $R(t)$ of the process is modeled as:
\begin{equation}
R(t) = R_{\text{init}} + \sum_{l=1}^{L} W_{l}
\end{equation}
where $R_{\text{init}}$ is the baseline framework memory footprint and $W_{l}$ is the physical size of layer $l$. For the 1.7B parameter Qwen3-TTS, the total weight size $W = \sum W_l \approx 3.8 \text{ GB}$, easily triggering Out-of-Memory (OOM) errors or heavy disk swapping on an 8GB laptop.

Under our proposed virtual address mapping, the file $F$ is bound to the virtual address space segment $\Omega_{\text{virtual}}$ via OS file descriptors, bypassing heap allocation. When layer $l$'s forward operator $f_l(\mathbf{x}_l; \mathbf{w}_l)$ executes, accessing the virtual address $\mathbf{w}_l \in \Omega_{\text{virtual}}$ triggers a Page Fault:
\begin{equation}
\text{PageFault}(\mathbf{w}_l) \implies \text{Kernel Page Cache Load}(F[offset_l : offset_l + size_l]) \to \Omega_{\text{physical}}
\end{equation}
Once the computation of layer $l$ completes, the operating system's LRU paging policy decommits the page, keeping the active physical RAM overhead $\Delta R(t)$ bounded by the size of the single largest layer:
\begin{equation}
\Delta R(t) \le \max_{l} (W_l) \ll \sum_{j=1}^{L} W_j
\end{equation}

\subsection{BFloat16 Numerical Compression}
\textbf{BFloat16} (Brain Float) compresses standard 32-bit float values ($\mathtt{Float32}$) down to 16 bits by keeping the exact same 8-bit exponent size, preventing dynamic range loss while cutting the parameter storage and transfer bandwidth by exactly 50\%. Mathematically, a BFloat16 value $v_{\text{bf16}}$ is represented as:
\begin{equation}
v_{\text{bf16}} = (-1)^{s} \times 2^{e - 127} \times \left(1 + \frac{m}{128}\right)
\end{equation}
where $s$ is the sign bit, $e$ is the 8-bit exponent, and $m$ is the 7-bit mantissa.

\subsection{Controlling Context-Switching Thrashing}
On client hardware with $P$ physical CPU cores, launching $T$ threads such that $T \gg P$ causes thread scheduler congestion. The execution time $T_{\text{exec}}$ is modeled as:
\begin{equation}
T_{\text{exec}} = \frac{T_{\text{calc}}}{k \cdot \min(T, P)} + \alpha \cdot \max(0, T - P)^{2}
\end{equation}
By forcing a strict thread limit of $T = \min(P, 4)$ using `torch.set_num_threads(4)`, we eliminate CPU core thrashing, stabilizing execution times.

\section{Benchmark Analysis and Results}
Comparative benchmarks were executed under unbuffered conditions to trace loading speed and memory overhead.

\begin{table}[h]
\centering
\caption{Qwen3-TTS CPU Performance Benchmarks}
\resizebox{\columnwidth}{!}{%
\begin{tabular}{lcc}
\toprule
\textbf{Evaluation Metric} & \textbf{Default Eager (Float32)} & \textbf{Optimized (BFloat16)} \\
\midrule
Model Loading Speed & 25.76 seconds & \textbf{8.28 seconds} \\
Initial RAM Overhead & 1,727.28 MB & \textbf{-428.54 MB} \\
Peak Inference RSS & ~2,108.51 MB & \textbf{~1,255.32 MB} \\
\bottomrule
\end{tabular}%
}
\end{table}

The most striking result is the \textbf{negative RAM overhead (-428.54 MB)} during optimized initialization. Because virtual memory mapping registers the parameter pages directly into the OS page tables, no eager heap memory allocation occurs. This allows active garbage collection to reclaim existing python memory overhead, bringing the net loading RAM footprint to virtual zero.

The generated output file has an audio length of 7.04 seconds at a 24,000 Hz sample rate. The exact file size is 337,964 Bytes. We prove the mathematical integrity of the BFloat16 synthesized output via the wave structural equation:
\begin{equation}
\text{Total Samples} = 7.04 \text{ s} \times 24,000 \text{ Hz} = 168,960
\end{equation}
\begin{equation}
\text{Audio Data Bytes} = 168,960 \times 2 = 337,920 \text{ B}
\end{equation}
\begin{equation}
\text{WAV File Size} = 337,920 + 44 \text{ B (Header)} = 337,964 \text{ B}
\end{equation}

\section{Theoretical Alignment: Weight Streaming and 15-bit Hangeul Compression}
The theoretical foundation of compressing parameter bandwidth is isomorphic to the phonological structure of Hangeul. Instead of representing Korean characters using a 3-byte (24-bit) UTF-8 structure, a syllable $S$ can be decomposed into Choseong (C), Jungseong (V), and Jongseong (T) phonemes, requiring exactly 5 bits each:
\begin{equation}
S_{\text{packed}} = (C \ll 10) \mid (V \ll 5) \mid T
\end{equation}
This compresses any Hangeul character into exactly 15 bits (a 37.5\% compression). Combining this bit-packed representation with virtual memory mapping enables ultra-fast, register-level prefix matching (e.g., initial consonant search via `S\_packed \& 0x7C00`) directly on mapped text corpora, bypassing heavy string parsing.

\section{Conclusion and Recommendations}
This research successfully validated and deployed virtual address space mapping, BFloat16 numerical compression, and thread limit scheduling on a real-world 1.7B parameter speech generative AI model. By restricting physical RAM load to virtual zero, we pave the way for democratized local AI deployment, allowing low-spec, legacy machines to run state-of-the-art models efficiently.

\end{document}
